\documentclass[11pt]{article}
\usepackage[margin=1in]{geometry}
\usepackage[T1]{fontenc}
\usepackage[utf8]{inputenc}
\usepackage{lmodern}
\usepackage{setspace}
\usepackage{hyperref}
\usepackage{booktabs}
\usepackage{enumitem}
\usepackage{graphicx}
\usepackage{amsmath,amssymb}
\usepackage[numbers]{natbib}
\usepackage{array}
\usepackage{longtable}
\usepackage{xcolor}
\onehalfspacing
\hypersetup{colorlinks=true,linkcolor=black,citecolor=black,urlcolor=blue}

\title{From Robot Tax to AI Fiscal Architecture:\\ \large Reframing Swiss Social-Insurance Financing in the Age of Generative AI}
\author{Graz Network}
\date{April 2026}

\begin{document}
\maketitle

\begin{abstract}
This article revisits a 2019 Swiss bachelor thesis that asked whether taxing robots is a desirable response to automation, labour-market disruption and pressure on social-insurance financing. The original thesis argued that robot taxation is generally not desirable in a small open economy because it generates high administrative costs, induces distortionary incidence effects, and rests on a weakly defined taxable object. I show that the core intuition of that thesis survives, but only after a major conceptual translation. Between 2019 and 2026, the frontier of automation shifted from embodied industrial robots toward generative AI, software agents, model APIs, enterprise copilots and other forms of intangible automation. At the same time, legal governance evolved toward risk-based obligations on firms and deployers rather than legal personhood for AI systems. Using a feasibility-adjusted framework for tax-base selection, the paper argues that the case against a literal robot tax is stronger in the age of generative AI than it was under the earlier Industry 4.0 framing. The more coherent policy response for a small open economy such as Switzerland is a fiscal architecture built around internationally coordinated firm-level taxation, selective broadening of social-insurance financing beyond payroll alone, active labour-market transition policy and operator-based AI regulation. The paper therefore moves the debate from ``Should robots be taxed?'' to ``Which tax bases remain administratively feasible, legally coherent and distributively adequate when automation is increasingly intangible, cross-border and embedded in corporate organization?''
\end{abstract}

\noindent\textbf{Keywords:} robots, artificial intelligence, automation, taxation, labour share, social insurance, Switzerland

\section{Introduction}
The proposal to tax robots has long served as a shorthand for a broader anxiety: if automation replaces labour, will tax and social-insurance systems built around wages lose both revenue and legitimacy? A 2019 bachelor thesis by David Graz approached this issue in the Swiss context and reached a cautious but clear answer: a tax on robots is generally undesirable in a small open economy because it is distortionary, difficult to administer and unable to guarantee an efficient redistributive outcome. The thesis nevertheless took seriously the underlying policy problem of declining labour share, rising inequality, and the financing of social insurance under automation.

That baseline question has become more urgent, not less. Yet the relevant technological object has changed. In 2019, the debate could still plausibly centre on ``robots'' as quasi-identifiable production assets. By 2026, much of the economically important automation frontier is driven not by physically embodied machines alone but by general-purpose AI models, software agents, workflow automation, model licensing, cloud-based inference and organizational systems that augment or substitute labour through intangible capital. Meanwhile, the regulatory environment has also changed. The European Union's AI Act entered into force on 1 August 2024 and is applying in phases through 2027. Switzerland decided in February 2025 to pursue a sector-specific AI-regulatory strategy and to prepare a consultation draft by the end of 2026. International tax coordination has also advanced through the OECD/G20 Pillar Two global minimum tax.

These developments matter because they change the terms of the original problem. The question is no longer simply whether a robot can be treated as a taxable object. It is whether any object-specific automation tax remains coherent once the production technology becomes distributed, intangible and cross-border. In this setting, the original thesis can be extended into a broader argument: the more automation migrates from identifiable machines to AI-enabled systems and intangible organizational capital, the weaker the case for taxing the object of automation and the stronger the case for taxing the firm, the rents it captures, or broader value bases that remain observable.

This article therefore makes three contributions. First, it reconstructs the economic logic of the original 2019 thesis and shows why its scepticism about a robot tax remains analytically valuable. Second, it explains how the shift from Industry 4.0 to generative AI intensifies the administrative and territorial weaknesses of object-based automation taxes. Third, it proposes a feasibility-adjusted framework for tax-base selection and derives from it a Swiss-oriented fiscal architecture centred on firm-level taxation, social-insurance base broadening, labour-market adjustment policy and operator-based AI regulation rather than AI legal personhood.

\section{Literature Review}
The classic robot-tax debate combined several strands of literature. One strand emphasized automation's capacity to raise productivity and expand welfare. Another highlighted labour displacement, polarization and the possibility that automation could weaken wage growth and labour's share of income. A third asked whether tax policy should slow or redirect the transition.

The original thesis drew on this tripartite debate. It engaged Daron Acemoglu and Pascual Restrepo's work on automation, task reallocation and the importance of new task creation. It also relied on Joao Guerreiro, Sergio Rebelo and Pedro Teles, who explored conditions under which taxing robots could reduce inequality but at a cost in efficiency. Uwe Thuemmel similarly found that gains from robot taxation are limited and highly sensitive to labour-market structure and policy design. Xavier Oberson, by contrast, explored the legal route of an ``electronic ability to pay,'' suggesting that if robots become sufficiently autonomous, some form of legal-fiscal personality might eventually be imagined.

From the standpoint of 2026, these sources remain central but must be re-read in light of two additional literatures. The first concerns generative AI, AI diffusion and productivity. OECD work now emphasizes that recent generative AI tools can improve performance on some tasks substantially, but that aggregate gains depend on organizational adoption, skills and institutional context. The second concerns the fiscal-policy response to AI. IMF work argues that the transition requires stronger social protection, investment in skills and tax systems that support broad sharing of the gains from AI rather than narrowly penalizing adoption.

The literature therefore points to a broad convergence. The concern that automation may intensify inequality has not disappeared. But the preferred policy response is increasingly moving away from blunt taxes on identifiable machines and toward the redesign of social-protection systems, coordinated corporate taxation, and policies that shape adjustment costs.

\section{From 2019 Robots to 2026 AI Systems}
The first reason the original thesis needs translation is technological. In the Industry 4.0 framing, a robot tax could be imagined as a tax on physical machinery, on purchases of robots, or on a firm's use of identifiable capital that substitutes for labour. Even then, the thesis correctly noted that software robots, location problems and legal ambiguity made administration difficult.

By 2026, however, the centre of gravity has shifted further toward intangible automation. Generative AI systems are delivered through APIs, enterprise suites and cloud infrastructure. Value creation is increasingly distributed across model providers, deployers, workflow designers, chip and compute providers, and multinational corporate structures. In many cases the ``robot'' is no longer a discrete object but a bundle of services, embedded capabilities and organizational redesign. A single firm can automate tasks without purchasing anything that looks like a taxable robot in the ordinary sense.

This matters for tax design because a tax must identify a base. Physical capital is never perfectly simple, but it is at least tangible. Diffuse AI capability is harder to locate and classify. Is the relevant taxable object the model? The inference call? The fine-tuned application? The subscription contract? The productivity gain? The saved headcount? The answer is not self-evident. Once the taxable object becomes ambiguous, administrative cost rises and behavioural avoidance becomes easier.

The second reason the thesis needs translation is legal. The AI Act regulates AI through a risk-based framework applied to providers, deployers, importers and distributors. The legal burden falls on firms and institutions, not on AI systems as rights-bearing or tax-bearing persons. Switzerland's chosen approach is similarly cautious and sector-specific. These developments suggest that real-world governance is converging on accountability-by-operator rather than personhood-by-machine.

The third reason is institutional. The original thesis worked in a national small-open-economy setting, which made relocation and competitiveness central. That remains true, but with broader channels. The mobility margin is no longer just the relocation of factories or machinery; it includes IP location, licensing structures, cloud contracting, and cross-border service delivery. Intangible automation therefore amplifies, rather than reduces, the small-open-economy problem identified in 2019.

\section{A Feasibility-Adjusted Framework}
To extend the thesis into a 2026 article, I propose a simple framework for comparing tax instruments. Let the policy desirability of tax base $i$ be:
\[
D_i = R_i - A_i - E_i - M_i,
\]
where $R_i$ denotes revenue and redistribution capacity, $A_i$ administrative and compliance cost, $E_i$ distortion cost, and $M_i$ the mismatch between the legal tax base and the actual locus of AI-enabled value creation.

This formulation is intentionally simple. Its purpose is not to yield point estimates but to organize the argument more rigorously than a purely verbal discussion.

\subsection{Revenue and redistribution capacity}
A tax instrument must raise revenue in a way that supports the redistributive or insurance objective that motivates intervention. A robot tax may appear attractive here because it symbolically targets capital that substitutes for labour. But symbolic targeting is not enough. If the base is narrow, avoidable or unstable, actual revenue may be weak.

\subsection{Administrative and compliance cost}
The 2019 thesis already emphasized the problem of definition, control and legal clarity. Under generative AI these costs rise. The more automation is provided through subscriptions, cloud services and integrated software suites, the more difficult it becomes to identify a taxable object without creating arbitrary thresholds or opportunities for relabelling.

\subsection{Distortion cost}
A tax on production inputs can alter firm behaviour in inefficient ways. In a small open economy, incidence may be transferred to labour, investment or location decisions. The original thesis captured this point well. Under today's conditions, distortion may take the form of delayed adoption, restructuring around tax arbitrage, or migration of contracts and IP rather than the obvious non-purchase of a machine.

\subsection{Mismatch between tax base and value creation}
This final term is what distinguishes the AI-era extension from the original thesis. A tax base can be observable yet conceptually misaligned with where value is created. For instance, taxing a robot purchase may miss value generated by AI licensing, model integration or data-driven coordination. As automation becomes intangible, the mismatch term becomes central.

\section{Legal and Fiscal Analysis}
This section compares four broad policy families.

\subsection{Object-specific robot or AI taxes}
These taxes target the automation object itself: the robot, the software system, the AI service or a proxy for machine substitution. Their intuitive appeal is political simplicity. They appear to restore symmetry between labour, which carries payroll and income-tax burdens, and automation capital, which seems to escape them.

Yet this family performs poorly in the feasibility-adjusted framework. First, definitional problems are severe. Second, location and sourcing are difficult in cross-border service environments. Third, firms can repackage functionality in ways that undermine the base. Fourth, the tax may penalize adoption in situations where AI is productivity-enhancing and complementary to labour rather than substitutive.

The older debate about granting AI systems legal personality was, in part, an attempt to solve the definitional problem by making the machine itself a legal bearer of obligations. But the policy trajectory in Europe and Switzerland has not moved in that direction. Real governance is being built around duties on firms, providers and deployers. For the foreseeable horizon, this makes AI personhood an analytically interesting but institutionally remote endpoint.

\subsection{Firm-level profit taxation and coordinated minimum taxation}
A second family targets the firm rather than the automation object. The strongest recent development here is Pillar Two. Switzerland introduced the OECD minimum tax from 1 January 2024 via a supplementary tax and extended implementation of the income inclusion rule from 1 January 2025. At the OECD level, the global minimum tax continues to be developed through consolidated commentary and new guidance.

This approach has several advantages. It taxes profit where the tax system already has concepts, enforcement structures and international coordination mechanisms. It is also more agnostic about the technological means by which profits are generated. If a firm captures rents from AI-enabled scale, pricing power or reorganization, profit-based instruments can reach those rents without needing to define a taxable robot.

This does not solve everything. Profit taxation remains vulnerable to international competition, design complexity and the limits of current allocation rules. But compared with an object-based robot tax, it fits the AI era better because it follows the firm and its income rather than trying to isolate a machine-like tax object.

\subsection{Broadening the social-insurance base}
The original thesis discussed alternatives such as gross value added and consumption-based finance. These options deserve renewed attention. If the core policy concern is that payroll-linked social-insurance financing becomes less reliable when production relies more on capital, software and AI, then one response is to rebalance the contribution base.

A broader value-added base has the advantage of not requiring the state to decide which machine or software system ``counts'' as a robot. Instead, it asks firms to contribute from a broader slice of value creation. This can improve neutrality between labour and capital. However, such broadening may also discourage investment if not designed carefully, especially in high-productivity sectors.

Consumption-based finance, including VAT-supported funding, has a different logic. It shifts financing toward final demand rather than production inputs. That may reduce direct penalties on adoption, but it raises distributional concerns because consumption taxes tend to be regressive unless combined with offsetting transfers or targeted support. The key lesson is that these instruments are imperfect, but they are still more administratively coherent than a literal robot tax under current technological conditions.

\subsection{Transition spending and labour-market policy}
A final family does not try to tax automation specifically at all. Instead, it focuses on financing adjustment: reskilling, active labour-market policies, wage subsidies, portable benefits and targeted support for displaced workers. IMF work strongly supports this direction. So does the World Economic Forum's repeated finding that skills gaps are a central barrier to successful transformation.

This is not a substitute for taxation, but it affects what taxation should do. The tax system may be better used to mobilize broad and stable resources than to target the automation object directly. In other words, the fiscal system should fund adaptation rather than attempt to micromanage the technological margin through a poorly targeted levy.

\section{Macroeconomic and Distributional Implications}
The 2019 thesis argued that in a small open economy, firms faced with a tax on automation inputs would often shift the burden through employment reductions, lower competitiveness or altered investment. That logic still holds. Indeed, it arguably holds more strongly when automation is embedded in organizational systems whose costs can be reallocated across jurisdictions and contracts.

At the same time, the macroeconomic environment is not purely one of substitution. OECD evidence on generative AI suggests meaningful productivity gains at the task level, but the aggregate outcome depends on adoption patterns, worker capability and institutional support. The WEF's 2025 survey of employers points in the same direction: substantial job displacement is expected, but so is substantial job creation, producing a net positive effect at the global level if transition is managed well. This means policy cannot be built on a simple presumption that all AI replaces labour. Some uses will substitute, some will complement, and many will recompose tasks.

This ambiguity strengthens the case against a literal robot tax. If the state cannot easily distinguish welfare-enhancing complementarity from harmful displacement, taxing the automation object risks blunting beneficial adoption while missing the rents and inequalities that actually matter.

Distributionally, the central challenge is still the labour share and the financing of social insurance. If high-value firms can scale output with relatively less labour, payroll-linked systems may weaken even if aggregate productivity rises. That suggests a need to decouple part of social-insurance financing from wages alone. A broader tax base, coordinated corporate taxation and well-designed transfers are therefore more promising than an object-specific tax on AI systems.

\section{A Swiss-Oriented Policy Architecture}
For Switzerland, the most coherent response is not a robot tax but a four-pillar fiscal architecture.

\subsection{Pillar 1: preserve neutral firm-level taxation}
Switzerland should continue to prioritize internationally coordinated taxation of large corporate groups, including effective implementation of the OECD minimum-tax framework. This does not solve every distributional issue, but it is better aligned with AI-generated rents than a tax on notional robots.

\subsection{Pillar 2: broaden social-insurance financing selectively}
Where payroll alone becomes a weaker base, a partial broadening toward value added or other broader contribution bases should be considered. The design must be careful: the goal is not to punish productive investment but to reduce excessive dependence on labour income alone.

\subsection{Pillar 3: finance transition rather than penalize technology}
Public policy should invest in training, reallocation, wage support where needed, and institutional mechanisms that help workers move into complementary tasks. This is especially important because skills gaps have become a first-order transformation barrier.

\subsection{Pillar 4: regulate AI through firms and operators}
Legal governance should continue to attach duties to providers, deployers and institutions, not to fictional AI taxpayers. This is consistent with both the EU and Swiss regulatory trajectory. Transparency, accountability and sector-specific obligations are more realistic than electronic personhood.

\section{Discussion}
Two objections deserve attention. The first is political. A robot tax has rhetorical force because it appears to make automation ``pay its share.'' A broader profit- or value-based approach may look less direct. Yet policy should not confuse symbolic fit with administrative fit. The wrong tax base can produce the appearance of justice while delivering little revenue and many distortions.

The second objection is distributive. Some critics may argue that firm-level taxation is too indirect and that only a targeted levy on automation truly links the source of disruption to the source of funding. This concern is understandable, but it assumes that the source of disruption is readily observable. Under generative AI, that assumption is increasingly untenable. When the technological source is diffuse and service-based, the more reliable route is to tax the institution that captures the gains and to use those resources for transition and redistribution.

There are limits to this paper. It does not estimate the size of Swiss revenue under alternative reforms, nor does it model incidence quantitatively across sectors. Its contribution is conceptual and institutional. But that is also its value: it explains why the original robot-tax vocabulary has become less useful just as the underlying fiscal problem has become more important.

\section{Conclusion}
The central intuition of the 2019 thesis remains persuasive: taxing robots is generally not a desirable solution for a small open economy facing automation-driven change. By 2026, however, the argument should be stated more strongly. The rise of generative AI and intangible automation does not rescue the robot tax; it further undermines it. The taxable object is harder to define, harder to locate, and less closely tied to the actual locus of value creation than before.

The policy debate should therefore move beyond the question of whether robots should be taxed. The more relevant question is which tax bases remain administratively feasible, legally coherent and distributively adequate when automation is increasingly embedded in intangible capital, organizational systems and cross-border services. For Switzerland, the answer is a fiscal architecture that combines coordinated firm-level taxation, selective broadening of social-insurance financing, robust transition policy and operator-based AI regulation. That architecture is less dramatic than the idea of taxing robots, but it is better suited to the world that actually emerged.

\bibliographystyle{plainnat}
\begin{thebibliography}{99}
\bibitem{acemoglu2018}
Acemoglu, D. and Restrepo, P. (2018).
\newblock Artificial Intelligence, Automation and Work.
\newblock In \textit{Economics of Artificial Intelligence}.

\bibitem{bakom2025}
Federal Office of Communications (BAKOM) (2025).
\newblock Artificial Intelligence: Switzerland's regulatory approach, 12 February 2025.

\bibitem{efd2024}
Federal Department of Finance (2024).
\newblock Implementation of the OECD minimum tax rate in Switzerland, updated 4 September 2024.

\bibitem{euaiact2026}
European Commission (2026).
\newblock AI Act: application timeline and FAQs on implementation.

\bibitem{graz2019}
Graz, D. (2019).
\newblock \textit{Une etude prospective sur l'impact de la robotisation de l'economie sur la fiscalite et le financement des assurances sociales}.
\newblock Bachelor thesis, Formation universitaire a distance, Suisse.

\bibitem{guerreiro2019}
Guerreiro, J., Rebelo, S. and Teles, P. (2019).
\newblock Should Robots Be Taxed?
\newblock \textit{NBER Working Paper} 23806.

\bibitem{imf2024}
Brollo, F., Dabla-Norris, E., de Mooij, R., Garcia-Macia, D., Hanappi, T., Liu, L. and Nguyen, A. D. M. (2024).
\newblock Broadening the Gains from Generative AI: The Role of Fiscal Policies.
\newblock IMF Staff Discussion Note 2024/002.

\bibitem{oecd2025sme}
OECD (2025).
\newblock \textit{Generative AI and the SME Workforce: New Survey Evidence}.
\newblock OECD Publishing, Paris.

\bibitem{oecd2026ai}
OECD (2026).
\newblock AI use by individuals surges across the OECD as adoption by firms continues to expand.
\newblock Announcement, 28 January 2026.

\bibitem{oecd2026pillar2}
OECD (2026).
\newblock Global Minimum Tax and January 2026 Side-by-Side package.

\bibitem{oberson2017}
Oberson, X. (2017).
\newblock Taxing Robots? From the Emergence of an Electronic Ability to Pay to a Tax on Robots or the Use of Robots.
\newblock \textit{World Tax Journal}.

\bibitem{thuemmel2018}
Thuemmel, U. (2018).
\newblock Optimal Taxation of Robots.
\newblock CESifo Working Paper 7317.

\bibitem{wef2025}
World Economic Forum (2025).
\newblock \textit{The Future of Jobs Report 2025}.

\end{thebibliography}

\end{document}
